\section{硬件案例深拆：从事件、对象到证据}

本节把硬件知识压到具体路径里。每个案例都回答五个问题：事件从哪里来，经过哪些硬件，内核维护哪些对象，用户看见什么，失败时从哪里取证。这样写的目的不是多堆页数，而是把“计算机相关硬件”变成能解释、能调试、能测试的系统知识。

\subsection{案例零：一条 C 语句如何落到硬件}

先看最普通的一句代码：

\begin{lstlisting}
array[i] = array[i] + 1;
\end{lstlisting}

编译器可能把它变成地址计算、load、加法、store。CPU 执行时并不是直接“改数组”，而是一步步经过寄存器、地址生成、TLB、页表权限、cache、store buffer 和一致性协议。

\begin{lstlisting}
1. compute virtual address: base + i * element_size
2. look up TLB for virtual page
3. if TLB misses, walk page table
4. check user/supervisor and write permission
5. load cache line into L1 if absent
6. read old value from cache line
7. add one in ALU
8. place store into store buffer
9. make modified cache line visible by coherence rules
\end{lstlisting}

这条路径解释了几个常见现象。若虚拟地址没有合法 VMA，内核给进程发送 \code{SIGSEGV}；若 VMA 合法但 PTE 不 present，内核可能分配页、读取文件页或执行 COW；若两个 CPU 更新同一 cache line 上的不同计数器，性能会因为伪共享下降；若 store 还在 store buffer 中，另一个 CPU 看到的顺序可能需要原子操作和内存屏障来约束。

\begin{longtable}{p{0.2\linewidth}p{0.35\linewidth}p{0.33\linewidth}}
\toprule
硬件步骤 & 出错或变慢的原因 & OS 证据 \\
\midrule
地址计算 & 指针越界、整数溢出 & sanitizer、异常地址、core dump \\
TLB/page walk & 页表不存在、权限不够 & page fault 统计、\code{perf stat} TLB miss \\
cache 访问 & cache miss、伪共享 & \code{perf stat} cache miss、LLC miss \\
store 可见性 & 数据竞争、缺少屏障 & KCSAN、TSAN、lockdep、代码审查 \\
\bottomrule
\end{longtable}

所以“内存访问”不是一个单点动作。操作系统中的缺页、COW、mmap、锁、原子变量、NUMA 和性能分析，都在为这条路径服务。

\subsection{案例一：复位到第一个用户进程}

从按下电源到第一个用户进程运行，机器经历了多级控制权转移。每一级都只拥有下一步所需的最小能力。

\begin{lstlisting}
power good
  -> reset vector
  -> immutable BootROM or firmware entry
  -> DRAM training and platform init
  -> bootloader or EFI stub
  -> kernel decompression and early entry
  -> early page tables and console
  -> memory allocator and scheduler
  -> driver probing and root filesystem
  -> exec first user process
\end{lstlisting}

这条链路里，硬件和软件界限不断变化。UEFI 固件初始化 DRAM、PCIe、ACPI 和基础平台状态；bootloader 或 EFI stub 把内核镜像、initrd、命令行、UEFI 内存地图传给内核；内核接管页表、中断、调度和驱动；最后 \code{execve} 把第一个用户态程序装入地址空间。

\begin{longtable}{p{0.2\linewidth}p{0.36\linewidth}p{0.32\linewidth}}
\toprule
阶段 & 最关键的硬件事实 & 失败证据 \\
\midrule
复位向量 & CPU 从固定地址取指 & 无串口输出、BMC 事件、逻辑分析仪 \\
DRAM 初始化 & 内存控制器和内存条必须稳定 & 内存训练失败码、固件日志 \\
固件表 & 内存地图、设备、中断、NUMA 信息 & ACPI 错误、\code{dmesg} 启动日志 \\
内核早期 & 栈、BSS、页表、early console & early panic、异常寄存器 \\
驱动 probe & 资源、IRQ、DMA、clock、reset & probe fail 日志、sysfs 设备和 ACPI 节点 \\
用户态启动 & rootfs、init、ELF、动态链接器 & kernel panic、\code{No init found} \\
\bottomrule
\end{longtable}

Linux 源码阅读可以从 \filepath{arch/x86/boot/}、\filepath{arch/x86/kernel/head64.c}、\filepath{init/main.c}、\filepath{drivers/firmware/efi/}、\filepath{drivers/acpi/}、\filepath{drivers/pci/} 开始。读的时候不要只看函数名，要问：这一行代码依赖哪些硬件已经可用，哪些还不可用。

\subsection{案例二：一次定时器抢占}

抢占式调度不是调度器单独完成的。它依赖硬件定时器在未来某个时间点打断 CPU。假设用户程序正在死循环，定时器中断到来后，CPU 切入内核入口，内核记录当前任务运行时间，设置重新调度标志，在返回路径切换到另一个任务。

\begin{lstlisting}
user loop running
  -> local APIC timer or TSC deadline timer fires
  -> CPU enters interrupt vector
  -> kernel saves interrupt frame
  -> tick handler updates scheduler clock
  -> current task exhausts slice
  -> return path calls scheduler
  -> switch registers and stack to next task
  -> return to user mode of next task
\end{lstlisting}

这里至少有四种硬件状态参与：定时器、当前 RIP 和寄存器、内核栈、页表或地址空间标识。若中断被长期关闭，抢占会延迟；若时钟源不可靠，调度统计会错；若保存寄存器不完整，任务恢复后会计算错误。

\begin{longtable}{p{0.2\linewidth}p{0.35\linewidth}p{0.33\linewidth}}
\toprule
内核对象 & 对应硬件状态 & 典型源码入口 \\
\midrule
\code{task\_struct} & 任务身份、调度状态、地址空间 & \filepath{include/linux/sched.h} \\
内核栈 & 中断和系统调用时的可信栈 & \filepath{arch/x86/entry/entry\_64.S} \\
调度类 & runnable 任务的选择策略 & \filepath{kernel/sched/} \\
clockevent & 产生下一次定时事件 & \filepath{kernel/time/} \\
上下文切换代码 & 保存恢复寄存器和栈 & \filepath{arch/x86/kernel/process\_64.c} \\
\bottomrule
\end{longtable}

可观测证据包括 \code{cat /proc/interrupts} 中 timer 相关计数、\code{perf sched}、\code{trace-cmd record -e sched:*}、\code{/proc/schedstat}。若系统卡顿，要区分是没有定时器中断、调度器没选到任务、任务不可运行，还是 CPU 被更高优先级中断或关中断区域占用。

\subsection{案例三：TLB shootdown 为什么会拖慢多核}

进程取消映射一段内存时，内核修改页表。但其他 CPU 可能仍然缓存旧 TLB 项。若内核立即释放物理页，另一个 CPU 可能继续用旧翻译访问已释放页，破坏隔离。因此内核必须让所有相关 CPU 失效旧翻译。

\begin{lstlisting}
unmap virtual page:
  take mm locks
  remove or modify PTE
  collect CPUs that may run this address space
  send IPI to those CPUs
  each CPU invalidates TLB entry or range
  wait for acknowledgement
  free physical page after old translations are gone
\end{lstlisting}

这就是为什么频繁 \code{munmap}、频繁改变页权限、频繁释放大量映射，在多核系统上可能很贵。成本不只是改内存里的页表项，还包括跨 CPU IPI、等待远端 CPU 响应、打断远端工作流和刷新 TLB 热状态。

\begin{longtable}{p{0.2\linewidth}p{0.36\linewidth}p{0.32\linewidth}}
\toprule
设计选择 & 好处 & 代价 \\
\midrule
立即 shootdown & 隔离清晰，回收及时 & IPI 和等待成本高 \\
批量 unmap & 分摊 shootdown 成本 & 回收延迟、实现复杂 \\
PCID & 降低地址空间切换刷新 & 标签管理和一致性仍复杂 \\
大页 & 减少 TLB 项数量 & 拆分和权限粒度成本 \\
\bottomrule
\end{longtable}

Linux 相关路径可读 \filepath{mm/mmu\_gather.c}、\filepath{arch/x86/mm/tlb.c}、\filepath{include/linux/mmu\_gather.h}。观察可以用 \code{perf stat -e tlb:tlb\_flush} 或架构支持的 TLB 事件，再配合 \code{perf record} 看是否大量时间落在 TLB flush 和 IPI 路径。

\subsection{案例四：一次 \code{malloc} 到物理页}

\code{malloc} 通常先在用户态分配器里找空闲块。若分配器没有合适空间，它可能通过 \code{brk} 扩展堆，或通过 \code{mmap} 建立新映射。即使系统调用成功，物理页也不一定立刻分配。第一次写入时触发缺页，内核才分配零页或物理页。

\begin{lstlisting}
malloc large object:
  userspace allocator decides mmap
  kernel creates VMA
  syscall returns virtual address
  user first writes address
  CPU raises page fault because PTE absent
  kernel checks VMA permissions
  allocator takes physical page from buddy/slab path
  kernel zero-fills page
  kernel installs writable PTE
  faulting instruction retries
\end{lstlisting}

这个案例把用户分配器、VMA、缺页、物理页分配、cache 和安全连接起来。零填充不是浪费，而是隔离要求：新进程不能读到旧进程释放页里的数据。NUMA first-touch 也在这里生效：哪个 CPU 第一次写页，内核通常倾向从本地 NUMA node 分配页。

\begin{longtable}{p{0.22\linewidth}p{0.34\linewidth}p{0.32\linewidth}}
\toprule
现象 & 硬件/内核原因 & 证据 \\
\midrule
\code{malloc} 快，首次访问慢 & 懒分配和缺页 & \code{perf stat} page-faults \\
跨 NUMA 变慢 & 页在远端 node & \code{numastat}、\code{numactl} \\
RSS 小于虚拟内存 & VMA 已建但页未实际分配 & \code{/proc/<pid>/smaps} \\
释放后内存不立刻回系统 & 分配器缓存和 VMA 策略 & \code{mallinfo}、\code{pmap} \\
\bottomrule
\end{longtable}

源码入口包括 \filepath{mm/mmap.c}、\filepath{mm/memory.c}、\filepath{mm/page\_alloc.c}、\filepath{mm/huge\_memory.c}。学习时要把“虚拟地址可用”和“物理页已经分配”分开。

\subsection{案例五：一次 DMA 描述符顺序错误}

假设驱动要让设备读取一段 buffer。正确顺序是：准备 buffer，映射 DMA，填写描述符，执行 DMA 写屏障，再写 doorbell。若屏障缺失，设备可能先看到 doorbell，再读到旧描述符或半初始化描述符。

\begin{lstlisting}
correct submit:
  desc.addr = dma_addr
  desc.len = len
  desc.flags = READ
  dma_wmb()
  writel(new_tail, DOORBELL)

buggy submit:
  desc.addr = dma_addr
  desc.len = len
  writel(new_tail, DOORBELL)
  desc.flags = READ
\end{lstlisting}

这个 bug 难复现，因为它依赖 CPU store buffer、cache、互联排序和设备读取时机。低负载、单核或某种 CPU 上可能一直正常；换平台、换优化级别、换设备后就出现随机超时或 DMA 错误。

\begin{longtable}{p{0.22\linewidth}p{0.34\linewidth}p{0.32\linewidth}}
\toprule
错误 & 后果 & 证据 \\
\midrule
doorbell 早于描述符可见 & 设备读到旧命令 & command timeout、设备错误状态 \\
completion 前释放 buffer & 迟到 DMA 写坏内存 & KASAN、IOMMU fault、随机崩溃 \\
方向标错 & cache clean/invalidate 错 & 数据旧、数据损坏 \\
unmap 漏掉 & IOVA 泄漏或权限残留 & DMA debug、IOMMU 统计 \\
\bottomrule
\end{longtable}

Linux 提供 DMA API 和 DMA debug 是为了让驱动不要直接拼物理地址。相关入口包括 \filepath{Documentation/core-api/dma-api.rst}、\filepath{kernel/dma/}、\filepath{lib/dma-debug.c}、\filepath{drivers/iommu/}。驱动评审时必须追踪 buffer 所有权：CPU 写、设备读、设备写、CPU 读，每次转换都要有对应屏障和 map/unmap。

\subsection{案例六：冷缓存文件读取}

用户调用 \code{read(fd, buf, len)}。若 page cache 没有对应页，路径会从 VFS 进入文件系统，再到块层、I/O 调度、NVMe 或其他存储驱动，最后设备 DMA 把数据放进页缓存。

\begin{lstlisting}
cold read:
  user read()
  fd -> file -> inode
  page cache lookup misses
  filesystem maps file offset to disk blocks
  block layer builds bio/request
  driver submits command to device queue
  device DMA fills page cache page
  completion interrupt wakes waiting task
  kernel copies data to user buffer
\end{lstlisting}

关键是区分三种“完成”。设备 completion 表示硬件完成了 DMA；page cache uptodate 表示内核页里有了数据；\code{read} 返回表示数据已经复制到用户 buffer。这三个时间点不同。

\begin{longtable}{p{0.2\linewidth}p{0.36\linewidth}p{0.32\linewidth}}
\toprule
层次 & 对象 & 源码入口 \\
\midrule
系统调用 & fd、file、iov iterator & \filepath{fs/read\_write.c} \\
页缓存 & page/folio、xarray & \filepath{mm/filemap.c} \\
文件系统 & inode、extent、block mapping & \filepath{fs/ext4/}、\filepath{fs/xfs/} \\
块层 & bio、request、blk-mq queue & \filepath{block/} \\
设备驱动 & submission/completion queue & \filepath{drivers/nvme/host/} \\
\bottomrule
\end{longtable}

证据可以用 \code{strace} 看系统调用，用 \code{vmtouch} 或 \code{mincore} 判断缓存，用 \code{blktrace} 或 \code{bpftrace} 观察块 I/O，用 \code{iostat -x} 看设备队列。若第二次读取明显更快，通常是 page cache 命中，不是硬盘突然变快。

\subsection{案例七：写文件、掉电和 \code{fsync}}

写文件最容易误解。应用调用 \code{write} 成功，只表示内核接收了数据，通常还在 page cache 中。后台 writeback 会稍后把 dirty page 写到设备。若应用需要崩溃后数据存在，必须使用 \code{fsync} 或文件系统提供的事务语义。

\begin{lstlisting}
safe replace pattern:
  write new contents to temp file
  fsync(temp file)
  rename temp file to target name
  fsync(parent directory)
\end{lstlisting}

这段模式看起来繁琐，是因为硬件和文件系统有多个持久化对象：文件数据、inode 元数据、目录项、日志提交、设备 volatile cache。只同步文件不一定同步目录项；设备返回写完成不一定意味着掉电后仍存在，除非 flush/FUA 语义把数据推到稳定介质。

\begin{longtable}{p{0.2\linewidth}p{0.36\linewidth}p{0.32\linewidth}}
\toprule
阶段 & 崩溃后风险 & 约束方式 \\
\midrule
\code{write} 返回 & 数据可能只在内存 & \code{fsync}、同步写策略 \\
writeback 完成 & 设备缓存可能易失 & flush、FUA、掉电保护 \\
日志提交 & 元数据恢复有序 & journaling、barrier \\
rename 完成 & 目录项可能未持久化 & fsync parent directory \\
\bottomrule
\end{longtable}

源码入口包括 \filepath{mm/page-writeback.c}、\filepath{fs/sync.c}、\filepath{fs/ext4/fsync.c}、\filepath{block/blk-flush.c}。测试应使用故障注入或虚拟块设备模拟掉电，而不是只在正常关机后看文件是否存在。

\subsection{案例八：USB 键盘按键到 shell}

USB 键盘不是把字符直接交给 shell。设备周期性上报 HID report，host controller 用 DMA 把报告放入内存，USB core 完成 URB，HID 驱动解析按键，input 子系统产生事件，tty line discipline 把按键转成终端输入，shell 的 \code{read} 被唤醒。

\begin{lstlisting}
key press:
  key matrix changes
  keyboard firmware builds HID report
  USB host polls interrupt endpoint
  DMA writes report to memory
  HID driver parses usage code
  input subsystem emits key event
  tty handles canonical mode and echo
  shell read() returns after line discipline policy
\end{lstlisting}

同一个按键可能因为键盘布局、修饰键、终端模式、输入法、tty 设置而得到不同字符。硬件只报告按键状态，字符语义在更高层解释。

\begin{longtable}{p{0.22\linewidth}p{0.34\linewidth}p{0.32\linewidth}}
\toprule
失败点 & 表现 & 证据 \\
\midrule
USB 枚举失败 & 设备不存在 & \code{lsusb}、\code{dmesg} \\
HID report 异常 & 按键错乱或无响应 & \code{hid-recorder}、内核日志 \\
input 层未创建设备 & 桌面无输入 & \code{/proc/bus/input/devices} \\
tty canonical 模式 & 输入到回车才返回 & \code{stty -a} \\
shell 被信号打断 & \code{read} 返回错误或重启 & \code{strace} \\
\bottomrule
\end{longtable}

相关源码入口有 \filepath{drivers/usb/core/}、\filepath{drivers/hid/}、\filepath{drivers/input/}、\filepath{drivers/tty/}。这个案例说明“输入”跨越硬件、固件、总线、驱动、内核子系统和用户态协议。

\subsection{案例九：网卡收包到 \code{recv}}

网卡收包路径是硬件队列和协议栈协作的经典例子。驱动预先把 RX buffer 映射给设备，设备收到包后 DMA 写入 buffer，写 completion，触发中断。内核通常用 NAPI 关闭一段时间中断并轮询多个包，减少中断风暴。

\begin{lstlisting}
packet receive:
  NIC receives bits from PHY
  MAC validates frame and writes DMA buffer
  NIC writes descriptor status
  MSI-X interrupt targets selected CPU
  NAPI poll harvests packets
  driver builds skb
  IP/TCP stack processes headers
  socket receive queue gets data
  epoll or recv wakes user task
\end{lstlisting}

网络路径里的硬件细节会直接影响上层延迟。RSS 决定流分到哪个队列；IRQ affinity 决定中断在哪个 CPU；NUMA 决定 DMA buffer 离 CPU 远近；interrupt moderation 决定中断合并延迟；GRO 决定包在协议栈里的聚合粒度。

\begin{longtable}{p{0.22\linewidth}p{0.34\linewidth}p{0.32\linewidth}}
\toprule
调优项 & 改变什么 & 风险 \\
\midrule
RX ring size & 可暂存包数量 & 太小丢包，太大延迟和内存占用 \\
interrupt moderation & 中断频率 & 吞吐和延迟折中 \\
RSS indirection & 流到队列映射 & 负载不均或破坏局部性 \\
IRQ affinity & completion 运行 CPU & 热点 CPU 或跨 NUMA \\
GRO/TSO offload & 包处理粒度 & 抓包和延迟理解更复杂 \\
\bottomrule
\end{longtable}

源码入口包括 \filepath{net/core/dev.c}、\filepath{net/ipv4/tcp\_input.c}、\filepath{drivers/net/ethernet/}。证据命令包括 \code{ethtool -S}、\code{ethtool -l}、\code{cat /proc/interrupts}、\code{nstat}、\code{ss -tin}、\code{dropwatch}。

\subsection{案例十：Wi-Fi 连接突然不稳定}

Wi-Fi 不稳定时，应用只能看到请求慢或断线。底层可能是射频信号差、漫游、信道拥塞、省电策略、固件崩溃、监管域限制或驱动队列堵塞。Wi-Fi 设备通常有固件，很多 MAC 层策略在固件中执行，OS 通过命令和事件与它协作。

\begin{lstlisting}
Wi-Fi path:
  supplicant requests connection
  driver sends commands to firmware
  firmware scans and authenticates
  radio transmits frames over selected channel
  firmware reports events and received frames
  kernel network stack exposes normal netdev
\end{lstlisting}

有线网络里 PHY/MAC 已经复杂，无线还多了信道、速率控制、重传、加密、漫游和省电。故障排查要看 \code{iw dev wlan0 link}、\code{iw dev wlan0 station dump}、\code{dmesg} 固件日志、\code{wpa\_supplicant} 日志、信号强度、重传计数和省电状态。

\subsection{案例十一：GPU 渲染一帧到屏幕}

用户界面的一帧通常由应用渲染、合成器合成、GPU 执行命令、显示控制器扫描输出。应用可能使用 OpenGL、Vulkan 或软件渲染；合成器把多个窗口合成到目标 buffer；KMS 在 vblank 边界 page flip。

\begin{lstlisting}
frame lifecycle:
  application renders into buffer
  compositor receives buffer and fence
  GPU executes composition commands
  fence signals render completion
  KMS schedules page flip at vblank
  display controller scans buffer
  monitor shows pixels
\end{lstlisting}

显示路径里有两个异步边界：GPU 命令完成和显示扫描完成。若 buffer 在 GPU 完成前被显示，会出现未完成内容；若 buffer 在显示扫描时被修改，会撕裂；若 page flip 错过 vblank，会多等一个刷新周期。

\begin{longtable}{p{0.22\linewidth}p{0.34\linewidth}p{0.32\linewidth}}
\toprule
对象 & 作用 & Linux 入口 \\
\midrule
buffer object & 显存或共享内存对象 & \filepath{drivers/gpu/drm/} \\
fence & 异步完成同步 & \filepath{drivers/dma-buf/} \\
plane/CRTC & 显示合成和扫描输出 & DRM/KMS helper \\
connector & HDMI、DP、eDP、DSI & EDID、hotplug、modeset \\
\bottomrule
\end{longtable}

黑屏或卡顿不能只怪应用。证据包括 \code{dmesg} GPU hang、\code{/sys/class/drm/} 连接器状态、compositor 日志、\code{modetest}、\code{intel\_gpu\_top} 或对应厂商工具。

\subsection{案例十二：音频爆音}

音频设备以固定采样率消耗数据。假设采样率 48kHz，period 只有几毫秒。用户态音频服务必须在设备读完当前 period 前填好后续数据。任何调度延迟、锁竞争、CPU 深睡眠、中断延迟都可能造成 underrun。

\begin{lstlisting}
audio underrun:
  device DMA consumes playback ring
  period interrupt should wake audio thread
  audio thread wakes too late
  next period contains no valid samples
  hardware outputs silence or repeats stale data
  user hears pop or gap
\end{lstlisting}

这个案例说明实时不是“CPU 总体很空”就够。即使平均 CPU 使用率只有 20\%，一次几十毫秒的调度停顿也能毁掉音频。证据包括 ALSA underrun 日志、音频服务日志、\code{cyclictest}、调度 trace、CPU C-state 和频率状态。

\subsection{案例十三：摄像头采集一帧}

摄像头路径从光开始。x86-64 PC 上常见的 USB UVC 摄像头会在设备内部完成传感器读出和图像处理，再通过 USB 主机控制器把帧数据传到内存；PCIe 采集卡则通过 BAR、MSI-X 和 DMA 队列交付帧。V4L2 或 media framework 把 buffer 交给用户态。每一帧都有时间戳和 buffer 生命周期。

\begin{lstlisting}
camera frame:
  sensor exposure
  USB xHCI or PCIe capture card transfers pixels
  host controller or device DMA writes output buffer
  MSI-X or completion event marks buffer done
  V4L2 dequeues buffer to userspace
  application encodes or displays frame
\end{lstlisting}

摄像头比普通字符设备复杂，因为它是 pipeline。sensor、CSI receiver、ISP、scaler、codec 都可能是不同子设备。格式、分辨率、帧率、lane 数、时钟和电源顺序必须一致。证据包括 \code{media-ctl -p}、\code{v4l2-ctl --all}、帧 drop 计数、CSI error 和 ISP 日志。

\subsection{案例十四：温度传感器触发降频}

CPU 或 GPU 高负载运行后温度升高。传感器超过 trip point，thermal framework 调用 cooling device，可能降低 CPU 频率、限制 GPU、提高风扇转速或最终关机。用户看到的是吞吐下降或延迟升高。

\begin{lstlisting}
thermal throttling:
  workload raises power
  temperature sensor crosses trip point
  thermal governor selects cooling state
  cpufreq or devfreq reduces frequency
  scheduler capacity model changes
  application observes lower throughput
\end{lstlisting}

硬件热状态是性能报告必须包含的上下文。证据包括 \code{/sys/class/thermal/thermal\_zone*/temp}、\code{cpupower frequency-info}、\code{turbostat}、\code{perf stat}、风扇转速和环境温度。若不记录这些，两个 benchmark 结果没有可比性。

\subsection{案例十五：Secure Boot 到磁盘解锁}

安全启动和磁盘加密常组合使用。固件验证 bootloader，bootloader 验证内核或测量内核，TPM PCR 记录启动链状态。磁盘密钥可以封存到特定 PCR 状态，只有启动链未被篡改时才自动解锁。

\begin{lstlisting}
measured boot:
  firmware measures bootloader into PCR
  bootloader measures kernel and initrd
  kernel or initramfs asks TPM to unseal key
  TPM checks PCR values
  if values match policy, releases secret
  disk mapper decrypts root filesystem
\end{lstlisting}

这条路径解决的是“离线篡改启动组件后窃取磁盘数据”的一部分风险。它不解决运行中 root 权限滥用，不解决内核漏洞，不解决用户已解锁后的数据泄露。安全硬件必须放在威胁模型中解释。

\begin{longtable}{p{0.22\linewidth}p{0.34\linewidth}p{0.32\linewidth}}
\toprule
机制 & 防护对象 & 边界 \\
\midrule
Secure Boot & 未授权启动组件 & 运行中漏洞仍可能攻破系统 \\
TPM sealing & 密钥只在匹配测量状态释放 & 授权状态下恶意软件仍可请求使用 \\
IOMMU & 设备 DMA 越权 & CPU 内核漏洞不由 IOMMU 防住 \\
内存加密 & 物理 DRAM 窃听或冷启动部分风险 & 运行中软件访问仍受权限控制 \\
\bottomrule
\end{longtable}

源码和工具入口包括 \filepath{drivers/char/tpm/}、\filepath{security/integrity/ima/}、\code{mokutil}、\code{tpm2\_pcrread}、\code{systemd-cryptenroll}。学习安全硬件时不要只问“开没开”，要问它保护哪条路径、失败时如何恢复。

\subsection{案例十六：虚拟机直通一块网卡}

设备直通把真实 PCIe 设备交给虚拟机使用。hypervisor 要从 host driver 手中解绑设备，绑定到 VFIO，建立 IOMMU domain，把设备 DMA 限制在 guest 内存范围内，并把中断重定向给 guest。

\begin{lstlisting}
device passthrough:
  host checks IOMMU group isolation
  host binds device to vfio-pci
  VM maps guest memory into IOMMU domain
  device DMA is translated by IOMMU
  MSI/MSI-X interrupts are remapped to guest
  guest driver thinks it owns a real PCIe device
\end{lstlisting}

直通性能好，但隔离依赖 IOMMU group、ACS、interrupt remapping 和设备本身可信度。若同一 group 中有多个不可隔离设备，把其中一个交给 guest 可能让 guest 影响 host 或其他设备。

\begin{longtable}{p{0.22\linewidth}p{0.34\linewidth}p{0.32\linewidth}}
\toprule
风险 & 原因 & 证据或约束 \\
\midrule
DMA 越界 & IOMMU 配置错误或关闭 & \code{dmesg} IOMMU、VFIO 日志 \\
中断注入错误 & interrupt remapping 问题 & VM 日志、host IRQ 统计 \\
group 不隔离 & PCIe 拓扑共享请求者 ID & \code{/sys/kernel/iommu\_groups} \\
设备 reset 不完整 & 设备状态泄漏到下个 owner & FLR 支持、驱动 reset 路径 \\
\bottomrule
\end{longtable}

相关入口包括 \filepath{drivers/vfio/}、\filepath{drivers/iommu/}、\filepath{drivers/pci/quirks.c}。这个案例把 PCIe、IOMMU、中断、安全和虚拟化连接在一起。

\subsection{案例十七：BMC 远程重启一台服务器}

服务器 BMC 可以在主机 OS 崩溃时仍然工作。管理员通过管理网络登录 BMC，查看传感器、日志和远程控制台，执行 power cycle。主 CPU 被断电重启，但 BMC 的固件和网络连接可能一直在线。

\begin{lstlisting}
BMC recovery path:
  administrator connects over management network
  BMC reads platform event log
  BMC controls power rails through board logic
  host CPU resets and firmware starts again
  OS later reads hardware error records
\end{lstlisting}

BMC 是运维救命工具，也是安全风险。它有独立凭据、独立网络栈、固件更新路径和高权限硬件控制能力。生产环境应隔离管理网，审计 BMC 账户，更新固件，记录远程操作。

\subsection{案例十八：JTAG、串口和早期调试}

当内核还没启动到日志系统，普通调试器和文件系统都不可用。x86-64 硬件调试依赖更低层入口：串口 early console、Intel DCI/JTAG、逻辑分析仪、BMC SOL、panic blink code、firmware debug port。

\begin{longtable}{p{0.2\linewidth}p{0.36\linewidth}p{0.32\linewidth}}
\toprule
调试入口 & 适用阶段 & 价值 \\
\midrule
UART early console & bootloader 和内核早期 & 最简单稳定的文本证据 \\
Intel DCI/JTAG & 裸机、固件和早期内核 & 停 CPU、看寄存器和内存 \\
BMC SOL & 服务器远程串口 & 无需现场接线 \\
logic analyzer & 信号级时序问题 & 查 I2C/SPI/GPIO/复位线 \\
crash dump & 内核运行后崩溃 & 保存内存和寄存器证据 \\
\bottomrule
\end{longtable}

这解释了为什么底层工程重视串口。屏幕和网络依赖太多驱动，串口依赖少，最适合定位启动早期、panic 和设备 probe 问题。

\subsection{硬件相关 Linux 源码阅读路线}

为了避免“工业对标只是概述”，这里给出一个实际阅读顺序。每个路径都建议先读数据结构，再读 probe，再读失败回滚，再读中断或 completion。

\begin{longtable}{p{0.25\linewidth}p{0.37\linewidth}p{0.26\linewidth}}
\toprule
主题 & 源码路径 & 先抓的对象 \\
\midrule
启动主线 & \filepath{init/main.c}、\filepath{arch/x86/kernel/head64.c} & early init、initcall \\
ACPI & \filepath{drivers/acpi/} & device 描述、电源状态和资源 \\
PCIe & \filepath{drivers/pci/} & \code{pci\_dev}、BAR、MSI \\
USB & \filepath{drivers/usb/core/} & \code{urb}、endpoint、hub event \\
DMA/IOMMU & \filepath{kernel/dma/}、\filepath{drivers/iommu/} & mapping、domain、fault \\
块层 & \filepath{block/} & \code{bio}、\code{request}、blk-mq \\
NVMe & \filepath{drivers/nvme/host/} & queue、command、completion \\
网络 & \filepath{net/core/dev.c}、\filepath{drivers/net/} & \code{sk\_buff}、NAPI、netdev \\
DRM/KMS & \filepath{drivers/gpu/drm/} & GEM、plane、CRTC、fence \\
音频 & \filepath{sound/core/}、\filepath{sound/pci/hda/} & PCM ring、codec、stream \\
摄像头 & \filepath{drivers/media/} & media graph、V4L2 buffer \\
热管理 & \filepath{drivers/thermal/} & thermal zone、cooling device \\
电源管理 & \filepath{drivers/base/power/} & runtime PM、system sleep \\
安全芯片 & \filepath{drivers/char/tpm/} & command、PCR、session \\
\bottomrule
\end{longtable}

读源码要带着硬件问题读。例如读 \code{pci\_driver.probe} 时，问：资源哪里来，MMIO 如何映射，IRQ 如何分配，DMA mask 如何设置，失败时如何释放。读 NAPI 时，问：为什么中断处理不直接把所有包处理完，poll budget 如何平衡延迟和吞吐。

\subsection{硬件测试报告模板}

下面这个模板适合本章学习后的最小测试。它不要求写代码，但要求把硬件结论落到证据。

\begin{lstlisting}
Hardware report:
  Machine:
    CPU model, core/thread count, cache levels, NUMA nodes
    memory size, memory type if visible, ECC status if visible
  Firmware:
    BIOS/UEFI version, ACPI evidence, Secure Boot status
  PCIe:
    list NVMe/GPU/NIC path, BAR, driver, IRQ/MSI-X, IOMMU group
  Storage:
    device model, queue count, scheduler, flush support, SMART summary
  Network:
    link speed, queues, offload, IRQ affinity, error counters
  USB/Input:
    topology, HID devices, hotplug evidence
  Display/Audio:
    connectors, modes, audio devices, underrun or error counters
  Power/Thermal:
    temperature sensors, frequency governor, cooling devices
  Security:
    TPM, Secure Boot, IOMMU, kernel lockdown if enabled
  Failure case:
    choose one failure and trace hardware -> kernel -> user symptom
\end{lstlisting}

报告的判断标准：每个结论都要有命令输出或日志证据；每个设备都要能说出驱动对象；每个性能结论都要说明硬件瓶颈可能在哪里；每个安全结论都要说明防护边界。

\subsection{常见误区集中纠正}

\begin{warningbox}
\begin{itemize}
\tightlist
\item “CPU 访问内存就是读数组”：实际还经过虚拟地址、TLB、页表权限、cache 和异常路径。
\item “\code{write} 成功就是落盘”：通常只是内核接收，持久化还需要 writeback、flush 和文件系统顺序。
\item “中断越多越实时”：过多中断会造成 CPU 浪费和抖动，很多高性能路径会合并中断或轮询。
\item “DMA 就是更快复制”：DMA 是设备直接访问内存，关键是所有权、映射、cache 和完成回收。
\item “设备插上驱动就该工作”：还需要供电、时钟、复位、总线枚举、固件描述、IRQ、DMA 和权限。
\item “虚拟机里的物理地址是真的物理地址”：guest physical 还要经过二级页表到 host physical。
\item “安全硬件打开就安全”：安全硬件只保护特定边界，不能替代内核和用户态安全。
\item “性能只看算法”：cache、TLB、NUMA、频率、温度、I/O 队列和中断亲和性都可能主导结果。
\end{itemize}
\end{warningbox}

\subsection{本节能力验收}

读完本节后，不要求记住所有芯片型号，但要能做到以下事情：

\begin{rubricbox}
\begin{itemize}
\tightlist
\item 能从一条 C 内存访问讲到 TLB、页表、cache、store buffer 和缺页。
\item 能从电源复位讲到固件、bootloader、内核 early init 和第一个用户进程。
\item 能解释定时器中断如何让内核抢回 CPU。
\item 能解释 TLB shootdown 为什么是多核虚拟内存的昂贵操作。
\item 能解释 \code{malloc} 成功和物理页真正分配之间的差别。
\item 能给 DMA 提交流程写出正确的所有权和屏障顺序。
\item 能完整讲出冷缓存文件读取、可靠文件写入、USB 键盘、网卡收包、GPU 显示、音频播放、摄像头采集的路径。
\item 能把 Wi-Fi 不稳、NVMe 超时、黑屏、音频爆音、热降频、设备不可见拆成硬件层、驱动层、内核对象层和用户可见层。
\item 能根据 Linux 源码路径找到至少一个硬件子系统的 probe、资源申请、completion 和错误恢复代码。
\item 能写出一份本机硬件测试报告，并给每个结论附上命令或日志证据。
\end{itemize}
\end{rubricbox}
